\headline{{\textbf{\Large\sffamily The Top Chop:}}\\
          {\textbf{\large\sffamily The Tree of the Month Contest}}}
\label{2025-09-contest}
\noindent
September brought a rich palette of evergreens to the benches of Twickenham Bonsai Club, as members explored this month's verdant theme: \textit{``Evergreen''}.
With dense foliage, ancient trunks, and deep green hues, the trees showcased the enduring strength and beauty of species that thrive all year round.
Whether refined over decades or still developing, each entry stood as a celebration of persistence and elegance.
Each member may enter one tree (novices may enter more; only the top score counts).
Judging is anonymous and peer\--based across trunk, foliage, pot, and theme: up to 12 points per judge.

\topic{Tree of the Month}
\noindent
This month's contest was a close\--run showcase of bold conifers and seasonal resilience.

\begin{enumerate}
    \item \textbf{Chris Durne}'s Juniper (\textit{Juniperus Communis}, 184 points).
    
    \item \textbf{Tony Ulatowski}'s Juniper (\textit{Juniperus Communis}, 167 points).

    \item \textbf{Graham Stevenson}'s Azalea (\textit{Rhododendron Japonicum}, 166 points).

    \item \textbf{David Fishwick}'s Cotoneaster (\textit{Cotoneaster Microphyllus}, 162 points).

    \item \textbf{Steve Palmer}'s Picea Glauca (\textit{Picea glauca}, 131 points).

    \item \textbf{Terry Wilson}'s Yew (\textit{Taxus Baccata}, 126 points).

    \item \textbf{Grant Kinnaird}'s Juniper (\textit{Juniperus Communis}, 111 points).

    \item \textbf{Brendan Gibbs}' Yew (\textit{Taxus Baccata}, 104 points).

    \item \textbf{Paul Owen}'s Japanese Black Pine (\textit{Pinus Thunbergii}, 66 points).
\end{enumerate}

\noindent
\textbf{Chris Durne} earns well\--deserved top honours this month with his refined \textit{Juniper}: a striking example of vitality and balance.

\topic{Novice Tree of the Month}
\noindent
While the main contest flourished, this month's Novice category saw no entries. 
We encourage our newer members to bring their trees forward next month and share their work: every tree adds to the learning and the celebration.

\vspace{-0.25em}
\begin{center}
$\cdots$
\end{center}
\vspace{-0.25em}

\noindent
This month's theme inspired a strong field of entries rooted in strength and longevity. 
From textured bark to lush, year\--round foliage, each tree embodied the theme with confidence and clarity.
Looking ahead, October's theme invites transformation with \textit{``Autumn Colours, Any Tree''}.
We look forward to fiery leaves, warm tones, and vivid contrasts. 
Until next time, keep your branches balanced and your passion growing!
\closearticle
\columnbreak
